\section{Locked Candidate-class Verification}
\label{app:branch-elimination}

The out-of-class comparison values were next subjected to every presently applicable locked constraint.  A comparison is classified by an already-fixed physical ledger or by a diagnostic whose candidate set and tolerance were fixed independently.  Residual ranking is not an equation-selection rule.

The audit contains 103 candidate--test findings.  Eighty-four lie outside physical diagnostic intervals, one has a nonzero exact-ledger defect, one is conditional, nine tests are nondiscriminating, and eight remain numerically nondiscriminated.  The classification is dominated by 61 electroweak coefficient triplets, 21 alternative single-monomial electromagnetic replacements, and the two nonminimal radial projectors $\Pi_{15}$ and $\Pi_{17}$.

\subsection{Compact Branches}

The fine-structure diagnostic places $\Pi_{15}$ and $\Pi_{17}$ outside its interval by many orders of magnitude.  At fixed $\Pi_{13}$, however, all three terminal comparison coefficients $\sigma_8=-1,0,+1$ remain numerically nondiscriminated.  The quasar observable
\begin{equation}
 \Delta v_{\rm high-low}
 =\operatorname{median}(\Delta v_{\rm est})_{\rm high\ L}
 -\operatorname{median}(\Delta v_{\rm est})_{\rm low\ L}
 \label{eq:pr35-quasar-nondiscrimination}
\end{equation}
contains no $\sigma_8$, $c_{\rm bc}$, or $\Pi$ dependence in its released-preprint form.  Each sign branch therefore makes the same negative-sign prediction. This successful PR-I stress test supports their common reduct but cannot select among its expansions.

The Kerr residual bound \cite{Kerr_1963} is also unchanged when only $\sigma_8$ is varied at fixed $\Pi_{13}$, because the radial spectral gap is unchanged.  The magnetic law
\begin{equation}
 B_{\rm geo}^{\rm PR}
 =\frac{(\hbar/e)(Z_A/c_{\rm bc})}{A_{\rm proj}^{(\theta)}}
 \label{eq:pr35-magnetic-nondiscrimination}
\end{equation}
does inherit the small branch dependence of $c_{\rm bc}$.  Its present use is an inverse area constraint, however, and the realized source area is not fixed independently.  It consequently does not select a terminal sign without a separate measurement or derivation of $A_{\rm proj}^{(\theta)}$.

\subsection{PR-II and PR-III Branches}

The vectorlike $t=0$ hypercharge family has a nonzero defect under the fixed nonzero electromagnetic charge ledger.  Complete generation replicas with $N_{\rm gen}\ne3$ do not match observed multiplicity once sheet multiplicity is identified with physical generation multiplicity, while every positive replication count satisfies the local-anomaly and Witten gates.  The normalized complete/nonredundant realization theorem supplies the exact selection.

None of the 21 tested comparison single-monomial replacements for $D_1+D_2$ lies inside the locked fine-structure gate.  This classifies that finite candidate family without making a statement about arbitrary higher-order combinations.  Of the 64 tested electroweak triplets, only
\begin{equation}
 (a,b,k)=(2,-1,1),\quad(2,-1,2),\quad(2,-1,3)
\end{equation}
pass both locked mass gates.  All three therefore pass the diagnostic gate. Residual ranking is not an admissibility rule, and the exact normalized charged-adjoint theorem rather than best fit selects $k=2$ in the canonical response class.

The stress audit therefore supplies substantial physical discrimination. Appendices
\ref{app:terminal-orientation-rigidity} and \ref{app:d3-coefficient-rigidity} now prove $\sigma_8=-1$ and the representation factor $k=2$ inside explicit boundary and normalized projected-adjoint classes.  The complete real-symmetric gauge-index coefficient classification at the fixed $D_3$ Lorentz/momentum kernel simultaneously shows that the three broad structural gates alone retain every
real $k$.  The remaining branches are resolved by the exact boundary and canonical-response closure theorems.  Those target-independent principles are defining conditions of the completed admissible class.  A physical-$D_3$ Hessian satisfying them is an equivalent microscopic representation and
cannot introduce another coefficient.  The full matrix and reproducible classifications are provided with the source archive.
